\documentclass[fleqn,10pt]{letter}
\usepackage[utf8]{inputenc}
\usepackage{units}
\usepackage{makecell}
\usepackage{float}
\usepackage[T1]{fontenc}
\usepackage[a4paper, margin=0.88in]{geometry}

\begin{document}

\begin{flushright}
Dr Carmen Cabrera\\
University of Liverpool\\
Liverpool, UK\\
29th of July, 2025\\
\end{flushright}

\begin{flushleft}
To Dr Steven Lukman \\
Chief Editor \\
\textit{Nature Communications} \\
\end{flushleft} 

Dear Dr Lukman,

I am writing on behalf of my co-authors, as we would greatly appreciate it if you could consider our manuscript titled "Sustained changes to urban mobility after COVID-19 amplified socio-economic inequalities in Latin America" for publication in \textit{Nature Communications}.

Our study presents a large-scale, longitudinal analysis of urban mobility recovery in Latin American countries following the initial COVID-19 disruptions. Using over 170 million anonymised mobile phone location records from Meta-Facebook users (March 2020–May 2022), we investigate how mobility patterns evolved over time, and how these changes intersect with socioeconomic and spatial inequalities.

Prior research has focused on short-term mobility declines in 2020, mostly in high-income countries. Our work fills a critical gap by extending the timeline through 2022, centring the analysis on cities in Latin America. We reveal that mobility recorded the sharpest declines and remained below pre-pandemic levels in most high-density and low socio-economic deprivation areas, while low-density and more deprived communities returned to baseline. These differences reflect the scale of the initial mobility shock, rather than subsequent recovery rates. 

Our findings illustrate how a global crisis can reinforce and prolong mobility-based inequalities, with important implications for economic resilience, social inclusion, and urban planning in unequal societies. We believe this work will be of broad interdisciplinary interest, particularly to researchers and policymakers focused on urban systems, mobility, and inequality.

We were encouraged to submit our manuscript to \textit{Nature Communications} by Dr William Burnside, following editorial feedback from \textit{Nature Cities}. We appreciate your consideration and look forward to the opportunity to share our findings with your readership.

%Our findings can help us understand how sustainable our cities are, which is an especially important issue given the rapid urbanisation process that the world is undergoing. Additionally, our results may be used to determine the top-priority regions to be targeted by policies for the alleviation of disruption caused by road accidents. We therefore want to make our work available to a broad, multidisciplinary audience and we believe that publishing in \textit{Nature Cities} can help us achieve this.

Finarlly, we would like to recommend three reviewers for our manuscript based on their expertise in human mobility, urban systems, and data-driven approaches, particularly through the use of location data derived from mobile phone devices:
\begin{itemize}
    \vspace{-1mm}
    \item Professor Esteban Moro; Northeastern University, College of Science
    \vspace{-1mm}
    \item Professor Ronaldo Menezes; University of Exeter, Department of Computer Science
    \vspace{-1mm}
    \item Dr Marta González; University of California, Berkeley, College of Environmental Design
\end{itemize}
\vspace{-1mm}
We thank you very much for the time invested on our behalf and look forward to hearing back from you.

Yours sincerely,

Dr Carmen Cabrera\\
Lecturer in Geographic Data Science\\
University of Liverpool

\end{document}